\subsubsection*{Equivariant realization of the selected row as cusp forms}
\label{subsec:selected-cusp-reading}

An index computed in geometric coordinates controls the selected row only through an equivariant identification. We exhibit the cusp-form realization on which the local cuts and the level actions agree.

Fix a lift $j$ and the run used at a coefficient place $v$.
Write $K_0$ for the construction field of the run, $K$ for
the common row field and $C_j$ for the rational cusp space
at its prescribed weight $\kappa$. The form coordinates
identify $C_j$ with a $K_0$-subspace of $K^{R_j}$ and
identify $K\otimes_{K_0} C_j$ with $K^{R_j}$. If \(b_1,\ldots,b_d\) is the common basis in the row data over \(K_0\), scalar extension preserves every coordinate independently:
\[
  K\otimes_{K_0}\bigoplus_{i=1}^dK_0 b_i=\bigoplus_{i=1}^dK(1\otimes b_i).
\]
 The chosen complex place determines an embedding
$\sigma_0:K_0\hookrightarrow\mathbb C$ and an extension
$\sigma:K\hookrightarrow\mathbb C$. Use the bijection
between the two-level components and the run's components.
At corresponding components the Hermitian data, frame
domains, embeddings and weights agree. Restrict the holomorphic forms to these frame domains.

The ordinary stable-range theta calculation in
Section~\ref{subsec:theta-row-rank} supplies the span before
projection. Here the target has rank $r$ over $E$ of
degree $2f$, the positive definite source has even rank
$n_V\ge 8fr+4$, and the level is odd with chosen split half
coprime to its conjugate.
The degree and rank conditions in the stable range combine to give
\[
  n_V\ge 8fr+4=4[E:\mathbb Q]r+4 .
\]
 Use the run's prescribed first
character Gram comparison and its admissible weight,
including positivity of the last weight. The chart
independence and the short theta transformation specialize
the stable unfolding to readings of raw short theta
sections. Every such short section is cuspidal at this
weight.

On each component extend a bounded holomorphic
section by zero off its frame domain, as a function on
all frame coordinates. The extension belongs to the
space on which the theta span is taken. The stable unfolding
expresses its family of components in a finite complex
span of the read short sections. Restricting the resulting
identity to the domains and transporting by the component
bijection gives the desired expansion of the original
family on the two-level model. The construction uses
the positive bad modulus and the Gram comparison of the
actual run.\label{cor:paired-row-reading-1}

\begin{proposition}[The equivariant cusp-form realization and the rank at the smaller level]
\label{prop:equivariant-selected-row}
For a lift \(j\) of either run in
Proposition~\ref{prop:linked-comparison-rows}, use its rational
cusp space \(C_j\) at weight \(\kappa\) over \(K_0\), with the
form coordinates of Definition~\ref{def:common-row-lattices}.
Extend along the complex embedding \(\sigma_0\) of \(K_0\). The chosen basis extends to a basis of the complex space of holomorphic cusp forms and identifies the corresponding character actions. The retained arithmetic projector is the same
one in the two realizations. For each lift the partition of the retained base labels gives the small-level image dimension
\[
 dZ_0(j)=D_0.
\]
 The full rank of the integral image uses the separate
off-prime trace inclusion.
\end{proposition}

\begin{proof}
\proofstep{Independence in the two coefficient realizations.}
For independence take a finite $K_0$-independent family of
cusp rules. Evaluate it on the framed test algebras through
all maps to $\overline{K_0}$. These evaluations detect a
relation of rules. Indeed the test algebras are reduced
and of finite type, so a nonzero test value survives at
a closed point over a finite extension of $K_0$. The span
of the evaluation rows is invariant under
$\operatorname{Gal}(\overline{K_0}/K_0)$, because composing
each test map with an automorphism retains it in the family.

If the row rank were deficient, the unique reduced row
echelon basis of the annihilator would be Galois invariant.
A nonzero basis row would give a $K_0$-relation among the rules. Hence some set \(I\) of rows, as large as the family of rules, gives a square matrix \(A_I\) with nonzero determinant. Under the chosen complex embedding \(\sigma\),
\[
  \det\big(\sigma(A_I)\big)=\sigma(\det A_I)\ne0 .
\]
 Framed base change realizes these tests
as complex frames. Uniformization and the common Borel
weight multiply all entries of an evaluation row by the
same nonzero scalar when passing to its frame-domain
representative. The resulting functions on frames are therefore independent. Restriction to the model domains preserves independence, since the functions vanish outside these domains and every component occurs.\label{cor:paired-row-reading-2}

\proofstep{The same rational character space on the two levels.}
Read the common $K_0$-basis in the two-level model and
extend its coefficients by $\sigma$. Independence and the
stable span give the isomorphism
\[
 \mathbb C\otimes_{K,\sigma} K^{R_j}\ \simeq\ V_\kappa .
\]
On a rational cusp form the lift of a level change acts on
the right in its moduli coordinate. In frame coordinates
this is the two-level pullback, with the same scalar and
the same change of component as in the run. The equality
on the basis therefore intertwines the full level action
after extension of scalars. In particular the dimensions
of the character kernel and its intersection with the kernel
of the good projector can be computed in $V_\kappa$.
Let \(A\) be the stacked character equations and \(P\) the good projector in these coordinates. Their common kernel is computed by the single matrix
\[
  \ker A\cap\ker P=\ker\begin{pmatrix}A\\P\end{pmatrix}.
\]
Changing the complex embedding preserves their algebraic
ranks.\label{cor:paired-row-reading-3}

\proofstep{Comparing with the retained arithmetic projector.}
The identification with the retained arithmetic rank uses
one more comparison at the selected point. In its row
coordinates the tie is the composite of the map to
$K\otimes_{K_0}C_j$ with the form-coordinate isomorphism
above. It identifies the subgroup of the arithmetic row
with the preimage of the base unipotent subgroup under the
level injection and identifies their characters there.
The determinant character in the coordinate change is
one on that subgroup. Consequently the arithmetic character
space maps bijectively to $V_\lambda$ after extension of
scalars.
Consequently the complex identification preserves the character rank, since scalar extension commutes with the common kernel
\[
  \bigcap_{u}\ker\big(\rho(u)-\mu(u)I\big).
\]
 The second row link uses the base-tuple theta
lift and its rational reading. In the split part of the
chosen realization these give the same reduced level
group and the same target projector, because the scalar
phase at the base is one.

The base-level lift identifies
the retained source space with
$V_\lambda\cap\operatorname{im}e_{\rm good}$ in the target
row. Its source dimension is the common $D_0$ of the
point, for every lift $j$. In the rank comparison the
base-tuple maps identify the dimensions of the summands
indexed by each retained irreducible type. Those summands
exhaust the respective spaces, so the sums of their
dimensions agree. Thus the target character rank
is the arithmetic $d_0$, and the retained rank is
\[
 D_0=d_0-\dim(V_\lambda\cap\ker e_{\rm good}).
\]
The equality uses invariance of the character space under
$e_{\rm good}$ and its decomposition into the image and
kernel of the restricted idempotent. In the integral
comparison all row lattices are transported by
this same tie. The scalar extension here computes ranks
and the lattice comparison uses localization in the
number field.\label{cor:paired-row-reading-4}

\proofstep{The rank in the determinant denominator.}
The small-level rank in the division denominator can be
read from the projection at the base tuple. The level
projection at the point's base index \(t_0\) is surjective
for each lift. The tuple readout sends \(t_0\) to \(\mu_0\),
since the base Hecke character is the identity and the
tuple characters distinguish the used indices.
On the matched row models the level projection induces the
theta adjoint \(\Theta_{\mu_0,j}^{*}\) by the class-sum identity
and the two level actions. Thus this adjoint surjects onto
the reference row space. Since

\[
  \ker\Theta_{\mu_0,j}=(\operatorname{im}\Theta_{\mu_0,j}^*)^\perp,
\]
\(\Theta_{\mu_0,j}\) is injective.

Decompose the reference row into the occurring class types.
The tied eigenspace multiplicity formula for each type gives
the sum of the row counts over that type's record labels.
These label sets partition the \(D_0\) base row labels,
and the class type summands form an internal direct sum.
Summing gives \(D_0\) at every lift. By injectivity the
small-level range \(\operatorname{im}\Theta_{\mu_0,j}\)
has that same dimension. This proves the identity
\(dZ_0(j)=D_0\) used in the cohomological division term.
The full rank of the projected \(J_v\) at a twisted index
follows separately from the local trace inclusion on the
assigned tuple.
\label{cor:paired-row-reading-5}
\end{proof}
